

Or ces
deux carrés sont les carrés «~de changement de base~» déduits des deux carrés (7.10.0), de sorte qu'on peut appliquer (IV 3). Il faut seulement vérifier que ces carrés correspondent à des couples @@M0@@ et @@M1@@ de schémas sur @@M2@@ qui sont tor-indépendants sur @@M3@@. Cela est trivial pour le premier couple, @@M4@@ étant plat, et aussi immédiat pour le deuxième, l'immersion @@M5@@ étant régulière de codimension 1 et restant ainsi après le changement de base @@M6@@.

Ne supposons plus maintenant @@M7@@ propre, mais supposons-le en revanche de tor-dimension finie. Je dis qu'il en est alors de même de @@M8@@ et @@M9@@. C'est trivial pour le premier, et pour le second résulte du fait (qu'on vient de signaler) que @@M10@@ et @@M11@@ sont tor-indépendants sur @@M12@@. On peut donc considérer le carré d'homomorphismes
@@M13@@
où les flèches horizontales sont définies en (IV 2.12). On notera que si @@M14@@ et @@M15@@ satisfont aux conditions de 7.7 i.e. sont à fibres régulières, alors (@@M16@@ étant régulier) @@M17@@ est automatiquement de tor-dimension finie, et le diagramme précédent peut s'interpréter comme la commutativité d'un diagramme d'anneaux @@M18@@.

Pour prouver la commutativité de (7.10.2), il suffit encore de prouver celle des carrés antérieurs du prisme analogue à @@M19@@, mais où les homomorphismes image directe @@M20@@, @@M21@@, @@M22@@ sont remplacés par les homomorphismes image inverse @@M23@@, @@M24@@, @@M25@@. Or la commutativité de ces carrés résulte aussi%
